% 第 7 章：功能模块
\section{功能模块}

\subsection{主界面交互}

\subsubsection{从电机角度到界面焦点}

主界面的关键不是单独实现一套滚动算法，而是把电机形成的物理档位直接解释为 LVGL 的焦点移动。FOC 任务持续读取轴角 $\theta_m$，并维护当前档位中心 $\theta_c$。主界面采用配置 1，档距约为 $8.23^\circ$，档位数量设为 0，表示菜单可以连续滚动而不设置机械意义上的首尾限位。当角度偏差
\begin{equation}
e_\theta=\theta_m-\theta_c
\end{equation}
超过档距与跨格阈值的乘积时，电机任务将 $\theta_c$ 平移一个档距，同时把逻辑位置变量 \texttt{position} 加一或减一。这样，连续角度先被触感状态机量化为稳定的整数档位；手停在吸附点附近时，位置值不会因传感器细小波动反复变化。

显示任务把旋钮注册为 \texttt{LV\_INDEV\_TYPE\_ENCODER}。输入回调 \texttt{encoder\_read()} 比较本次与上次的档位值：位置增大时输出 \texttt{enc\_diff=+1}，位置减小时输出 \texttt{enc\_diff=-1}。编码器输入与默认 \texttt{lv\_group} 绑定后，LVGL 按差分方向把焦点移动到前一个或后一个可交互按钮。因此，菜单选择形成了“轴角跨过档位阈值—档位整数变化—编码器差分—焦点切换”的逐级映射。第 6 章已经说明了 FOC 与 LVGL 分核运行及消息传递，本节关注的是这条映射如何把物理旋转转化为用户可见的选择结果。

\subsubsection{焦点的视觉表达与确认操作}

菜单容器使用纵向 Flex 布局，并开启纵向滚动和中心吸附。焦点变化后，目标按钮自动滚动到圆屏中央。为了让列表与圆形屏幕匹配，\texttt{scroll\_event\_cb()} 根据子项中心与容器中心的纵向距离 $d_y$ 计算横向偏移
\begin{equation}
x=r-\sqrt{r^2-d_y^2},\qquad |d_y|<r,
\end{equation}
超出半径时则把 $x$ 限制为 $r$。离中心越远的条目横向缩进越大、透明度越低，于是列表在视觉上呈现沿圆弧远离观察者的效果；位于中心的条目偏移最小且完全不透明，自然成为当前焦点。界面没有额外绘制固定光标，焦点位置、中心吸附和电机档位三者共同给出选择反馈。

GPIO0 的触摸输入在连续检测到有效触摸后向 LVGL 报告按下状态，当前焦点按钮随后产生 \texttt{LV\_EVENT\_CLICKED}。按钮回调负责创建目标页面、播放切屏动画，并按功能需要切换电机配置或工作模式。例如进入番茄钟时换为精细档位，进入音乐时发送播放事件，进入跳一跳时由游戏页面启动弹簧模式。旋转负责“选择”，按下负责“确认”，而每跨一格都能感到真实吸附，使视觉焦点与手上感知保持一致。

\subsection{物联网风扇控制}

\subsubsection{方案选择与通信实现}

这一模块控制的实机型号为 \texttt{xiaomi.fan.p69}。开始时我们考虑过米家云端、Home Assistant 中转和局域网直连三种方案。前两种方案分别需要处理账号登录或依赖一台常开的电脑，因此最后选择让 ESP32 通过家庭 Wi-Fi 直接与风扇通信。这样演示时只要旋钮和风扇处在同一网络中，就能独立完成控制。

米家风扇并不能直接接收普通字符串指令，所以我们在 \texttt{miio\_transport.cpp} 中补写了 miIO 通信过程，包括发现设备、整理请求内容、加密发送以及解析回复；在 \texttt{xiaomi\_fan\_p69.cpp} 中再把“读取状态、开关、调速、左转、右转”等操作封装成简单函数，供上层页面调用。最初参考资料给出的型号是 \texttt{p70}，但实机通信返回的是 \texttt{p69}。我们先用只读请求逐项核对电源、模式、风速和摆风状态，确认属性编号后才加入控制指令，避免因型号判断错误而反复修改界面代码。

\subsubsection{旋钮交互与实际调试}

风扇页面分为开关、风速和左右转动三个功能。旋转旋钮用于调节，触摸用于确认或返回；风速页用刻度盘和指针显示 $1\sim100$ 的数值，左右转动页则用圆弧表示操作方向。页面初版出现过中文字形缺失的问题，后来根据界面实际使用的文字重新生成字库，才保证了圆屏上的显示效果。

实机测试时还发现，如果每转过一格就立即等待风扇回复，画面会有明显停顿。为此我们把联网控制放到单独任务中：界面只记录用户最后调到的风速，网络任务按一定间隔发送；快速连续转动时，中间的旧数值会被新数值替代。左右转动指令也限制发送速度和累计数量，避免旋钮转得太快时风扇在随后很长一段时间内继续动作。这样处理后，电机手感和页面刷新不再被网络等待打断。

风速和转向页面都采用可连续旋转的齿轮手感，旋钮本身不设置硬限位，数值到达边界后由程序限制。每跨过一个档位，界面数值改变 2，既保留了清晰的档位感，也减少了从最低档调到最高档所需的圈数。目前实机已经完成状态读取、开关、风速和页面操作测试；左右转动指令已经能够发出，但单次动作对应的实际角度仍需继续根据风扇表现调整。

\subsection{番茄钟与音乐演奏}

\subsubsection{番茄钟的设计思路与核心功能}

番茄钟把计时与触觉提醒分开实现。页面提供工作、休息和退出三个操作，工作时间为 25 min，休息时间按演示配置为 2 min；FreeRTOS 软件定时器每秒更新一次 \texttt{MM:SS}，倒计时结束后启动提醒定时器。LVGL 仅负责按钮事件和时间显示。

进入页面时采用配置 2，即 256 个位置、$1^\circ$ 档距的精细档位，便于移动焦点。按钮确认和超时提醒均通过 \texttt{BUTTON\_CLICK} 触发 \texttt{motor\_shake()}，使电机产生短促振动；退出时停止定时器，并恢复主界面的连续档位配置。因此电机在操作阶段提供细腻的档位感，在计时结束后兼作触觉提醒器。

\subsubsection{电机发声原理}

无刷电机的绕组电流会产生磁场，并与永磁转子相互作用形成电磁转矩和径向力。当驱动电压按音频频率周期变化时，电磁力会激励转子、定子和外壳产生微小机械振动，结构表面推动空气形成声音。信号频率决定音高，持续时间决定节奏，电压、绕组阻抗和机械共振则影响响度与音色。

正常 FOC 根据转子角度合成旋转磁场，以获得平滑作用力；音乐模式则暂时停止 \texttt{loopFOC()}，主动向绕组施加周期激励，把机械振动作为主要输出。由于电机没有扬声器纸盆，频率响应也不平坦，因此其效果更接近蜂鸣器，旋律能够辨认但音色较为粗糙。

\subsubsection{音乐模式的实现}

音乐页面通过 \texttt{MUSIC\_PLAY} 通知电机任务接管播放。程序先令 FOC 输出为零，再把三个相线 PWM 引脚切换到 ESP32 LEDC 通道。旋律与时值保存在两个数组中，播放逻辑使用 \texttt{millis()} 非阻塞地推进音符；休止符输出零频率，其余音符按
\begin{equation}
f_{\mathrm{out}}=3f_{\mathrm{note}}
\end{equation}
提高到 3 倍后交给 \texttt{ledcWriteTone()}，从而按音符频率和时值组织旋律。

三个相线由 LEDC 产生随音符切换的周期性 PWM，使绕组电流形成相应频率的电磁激励，并通过电机结构振动输出可辨识的旋律。播放结束或触摸退出时，程序关闭音调与驱动并重启 ESP32，以重新初始化 SimpleFOC 和恢复旋钮手感。

\subsection{BLE 电脑控制}

\subsubsection{开源库的选择与接入}

BLE 部分采用了 BlynkGO 维护的开源项目 \href{https://github.com/BlynkGO/ESP32-BLE-Combo}{ESP32 BLE Combo Keyboard Mouse}。这个库可以让 ESP32 在同一次蓝牙连接中同时充当键盘和鼠标，并支持音量键等多媒体按键，正好覆盖本项目需要的滚轮、音量、组合键和鼠标按键功能。我们将库的源码放入 \texttt{lib/ESP32-BLE-Combo-main}，在 \texttt{bluetooth\_mouse.cpp} 中包含 \texttt{BleCombo.h}，通过 \texttt{Keyboard.begin()} 和 \texttt{Mouse.begin()} 启动蓝牙设备。之后调用 \texttt{Keyboard.write()}、\texttt{Keyboard.press()}、\texttt{Mouse.move()} 等接口即可向电脑发送标准 HID 操作，电脑端不需要再安装配套软件。

开源库解决的是蓝牙连接和按键报文发送，而旋钮怎样产生这些操作仍由我们自己完成。程序从 SimpleFOC 读取当前角度，以进入电脑模式或切换功能时的位置作为零点，再在 \texttt{run\_pc\_mouse\_logic()} 中把角度变化转换成键盘、鼠标命令，同时计算电机的回馈力度。这样，BleCombo 提供了电脑能够识别的输入接口，SuperKnob 原有的角度检测和力反馈则决定每一种操作的手感，二者在同一控制函数中结合起来。

实际运行时，BLE 与完整的 LVGL 界面同时启动会占用较多内存。我们因此把主界面和电脑控制分成两种启动模式：在菜单中选择“电脑”后，程序记录模式并软重启，只保留简化的 \texttt{PC MODE} 提示画面、FOC 和 BLE；按下硬件复位键则回到正常主界面。这个处理虽然需要一次短暂重启，但解决了早期蓝牙初始化不稳定的问题，也使演示过程更加可靠。

\subsubsection{四种控制方式与触感配合}

进入电脑模式后，触摸 GPIO0 可以依次切换表 \ref{tab:pc-control} 中的四种功能。每次切换都会把当前位置重新设为零点，并释放仍处于按下状态的 Alt 键或鼠标键，防止切换后电脑继续收到错误的长按输入。触摸程序必须检测到手指已经松开，才会接受下一次触摸，因此按住不放不会连续跳过多个模式。

\begin{table}[H]
\centering
\small
\begin{tabular}{p{0.17\linewidth}p{0.30\linewidth}p{0.39\linewidth}}
\toprule
功能 & 发送给电脑的操作 & 旋钮手感 \cr
\midrule
网页滚轮 & 每转过约 $0.20\ \mathrm{rad}$ 发送一次滚轮事件 & 一格对应一次滚动，具有清晰的棘轮感 \cr
音量调节 & 每转过约 $0.08\ \mathrm{rad}$ 发送一次音量增减键 & 中间区域顺滑，转到两端时逐渐出现阻力 \cr
应用切换 & 向一侧拨动后发送 Alt+Tab，继续转动可连续选择窗口 & 像带回弹的拨杆，回到中心时松开组合键 \cr
鼠标按键 & 向左或向右拨动，分别按下鼠标左键或右键 & 中心位置不触发，越过阈值后保持按下，松手回中后释放 \cr
\bottomrule
\end{tabular}
\caption{电脑控制的角度映射与触感设计}
\label{tab:pc-control}
\end{table}

滚轮和音量功能强调“转过一格、执行一次”，应用切换和鼠标按键则强调“拨动、保持、回中”。调试时我们分别调整了跨格角度、触发阈值和回中力度，使电脑收到的操作次数与手上的档位感尽量一致。最终这个模块并不是简单地把普通旋钮做成蓝牙按键，而是利用 SuperKnob 可改变手感的特点，让四种电脑操作在使用方式上也能被明显区分。

\subsection{钓鱼小游戏}

\subsubsection{把旋钮设计成阻尼卷线器}

钓鱼游戏没有沿用普通 \texttt{KnobConfig} 档位，而是把旋钮切换为纯粘性阻尼。进入游戏时页面调用 \texttt{update\_motor\_config(5)}，选中 \texttt{super\_knob\_configs} 中 \texttt{damping\_strength}=0.6 的配置 5；此时电机不再维护档位中心，而是输出与角速度成正比的反向力矩
\begin{equation}
\tau=-c\,\omega,
\end{equation}
其中 $\omega$ 为低通滤波后的角速度。该力矩没有回中项，旋钮转动时受到随速度增大的阻力，松手后停在哪就停在哪，像一根收线卷线器：转一下捕捉框移动一段，停止转动则框停在当前深度。

阻尼反馈直接使用传感器差分测速。12 位 AS5600 的量化噪声会使相邻采样产生约 $1.5\ \text{rad/s}$ 的速度尖峰，直接反馈会把噪声放大成持续抖动。为此固件先用一阶低通
\begin{equation}
\hat\omega_{k}=(1-\alpha)\hat\omega_{k-1}+\alpha\,\omega_k,\qquad \alpha=0.05
\end{equation}
磨平尖峰，再经过 $0.15\ \text{rad/s}$ 的静止死区，低于死区的残余波动直接输出零力矩。游戏页面的 \texttt{game\_tick()} 每 30 ms 读取一次旋钮相对转角 $\Delta\theta$，并按下式移动捕捉框
\begin{equation}
\Delta y=-\frac{\Delta\theta}{2\pi}\,H,
\end{equation}
其中 $H$ 为游戏区高度，即旋钮转满一圈正好把捕捉框从区底送到区顶。菜单阶段仍使用配置 1 的连续强档位，旋转有清晰的吸附感，供用户在开始、难度和最高分三个选项间切换。

\subsubsection{游戏状态与得分规则}

游戏开始后，代表鱼的圆点沿竖直方向随机游动。鱼在随机加速度之外还会以一定概率进入短时冲刺，速度限幅为 $\pm 3.5$ 倍难度系数，触到上下边界则反弹，因此玩家必须持续调整捕捉框去追踪鱼的位置。进度条按鱼与捕捉框的相对位置增减：当两者中心距离满足 $|y_{\text{fish}}-y_{\text{bar}}|\le h_{\text{bar}}/2+r_{\text{fish}}$ 时进度增加，否则下降；进度到达 100 记钓到一条，得分加一并通过 \texttt{BUTTON\_CLICK} 触发一次震动，随后放出一条新鱼并把进度重置为 50；进度归零则判定鱼已逃脱，页面记录本局成绩并在一段时间后自动回到菜单。

难度分为三档，在菜单中切换，对应鱼速倍率 $0.7$、$1.0$ 和 $1.5$。最高分保存在 \texttt{RTC\_NOINIT\_ATTR} 内存中，软重启不清零，仅彻底断电后才会复位，读取时还做了 $0\sim9999$ 的越界兜底。整个游戏形成“阻尼手感—捕捉框—鱼的位置—进度与得分”的闭环：旋钮的阻尼手感不再只是操作辅助，而是直接决定玩家能否稳定地把捕捉框停在目标深度上，手感本身就是游戏规则的一部分。

\subsection{跳一跳}

\subsubsection{把旋钮设计成蓄力弹簧}

跳一跳没有沿用普通 \texttt{KnobConfig} 档位，而是在电机任务中设置独立的 \texttt{MOTOR\_MODE\_JUMP}。开始一局时记录当前轴角为中心 $\theta_0$，旋转量为 $\Delta\theta=\theta_m-\theta_0$，电机输出采用弹簧与阻尼组合：
\begin{equation}
u_h=-3.2\Delta\theta-0.08\omega.
\end{equation}
比例项形成越拧越紧的回中力，速度项抑制回弹振荡；输出限制在 $\pm5.5$，超过 $100^\circ$ 后进一步增强反向作用力。蓄力比例为
\begin{equation}
c=\operatorname{clip}\left(\frac{|\Delta\theta|}{100^\circ},0,1\right)
\end{equation}
当峰值达到 $8^\circ$、角度回落至少 $3^\circ$ 且旋钮正向中心运动时，程序判定用户松手。FOC 任务把蓄力值和释放标志写入临界区保护的共享状态，LVGL 任务只读取结果。

\subsubsection{游戏状态与得分规则}

游戏定时器每 20 ms 更新蓄力条并检查释放事件。释放后，小方块沿抛物线起跳，距离、高度和动画时长均随蓄力比例增大：
\begin{align}
x_{\mathrm{end}}&=x_{\mathrm{start}}+175c,\\
h&=48+48c,\\
T&=500+280c\ \text{ms}
\end{align}
目标平台的位置和宽度随机生成。落在平台上可得分，落入中心区域还能获得连续加分；成功后生成下一平台，落空则进入下落和结束状态。该模块形成“旋转蓄力—电机回弹—释放检测—跳跃与落点反馈”的闭环，LVGL 只承担蓄力条、平台、方块和得分绘制，而弹簧手感本身就是游戏规则的一部分。
